% !TeX root = ../main.tex

% 符號列表。中括號裡是符號，大括號裡是說明。
% 下面每個類別各放一個例子，涵蓋數學符號、排版符號與英文縮寫三種標籤，
% 三者的排版方式相同。請留下與自己領域相關的幾項，其餘刪除；
% 整章不需要的話，把 main.tex 裡 % !TeX root = ../main.tex

% 符號列表。中括號裡是符號，大括號裡是說明。
% 下面每個類別各放一個例子，涵蓋數學符號、排版符號與英文縮寫三種標籤，
% 三者的排版方式相同。請留下與自己領域相關的幾項，其餘刪除；
% 整章不需要的話，把 main.tex 裡 % !TeX root = ../main.tex

% 符號列表。中括號裡是符號，大括號裡是說明。
% 下面每個類別各放一個例子，涵蓋數學符號、排版符號與英文縮寫三種標籤，
% 三者的排版方式相同。請留下與自己領域相關的幾項，其餘刪除；
% 整章不需要的話，把 main.tex 裡 % !TeX root = ../main.tex

% 符號列表。中括號裡是符號，大括號裡是說明。
% 下面每個類別各放一個例子，涵蓋數學符號、排版符號與英文縮寫三種標籤，
% 三者的排版方式相同。請留下與自己領域相關的幾項，其餘刪除；
% 整章不需要的話，把 main.tex 裡 \input{front/denotation} 那一行註解掉。
%
% 選用參數 [3cm] 是標籤欄寬度，預設 2.5cm。標籤比欄寬長時說明會被推開，
% 把數字調大即可。
%
% One example per broad field, covering the three kinds of label this
% environment handles: maths, typographic marks, and plain abbreviations. Keep
% the ones that fit your discipline and delete the rest. The optional [3cm] is
% the label column width; widen it if a label overruns.

\begin{denotation}[3cm]

% 1. 人文與文化 Humanities and culture
\item[$\dagger$]{
  obelus, marking a corrupt or doubtful passage in a critical edition
}

% 2. 社會、政治與公共政策 Social science, politics, and public policy
\item[GDP]{
  gross domestic product
}

% 3. 法律與制度 Law and institutions
\item[\S]{
  section sign, citing an article of a statute
}

% 4. 商業與管理 Business and management
\item[ROI]{
  return on investment
}

% 5. 基礎科學與數學 Basic science and mathematics
\item[$\hbar$]{
  reduced Planck constant
}

% 6. 工程、資訊與科技 Engineering, computing, and technology
\item[$O(\cdot)$]{
  asymptotic upper bound on a function's growth rate
}

% 7. 生物、農業與環境 Life science, agriculture, and environment
\item[$\mathrm{LD}_{50}$]{
  median lethal dose, the dose fatal to half of a test population
}

% 8. 醫學、健康與公衛 Medicine, health, and public health
\item[OR]{
  odds ratio
}

\end{denotation}
 那一行註解掉。
%
% 選用參數 [3cm] 是標籤欄寬度，預設 2.5cm。標籤比欄寬長時說明會被推開，
% 把數字調大即可。
%
% One example per broad field, covering the three kinds of label this
% environment handles: maths, typographic marks, and plain abbreviations. Keep
% the ones that fit your discipline and delete the rest. The optional [3cm] is
% the label column width; widen it if a label overruns.

\begin{denotation}[3cm]

% 1. 人文與文化 Humanities and culture
\item[$\dagger$]{
  obelus, marking a corrupt or doubtful passage in a critical edition
}

% 2. 社會、政治與公共政策 Social science, politics, and public policy
\item[GDP]{
  gross domestic product
}

% 3. 法律與制度 Law and institutions
\item[\S]{
  section sign, citing an article of a statute
}

% 4. 商業與管理 Business and management
\item[ROI]{
  return on investment
}

% 5. 基礎科學與數學 Basic science and mathematics
\item[$\hbar$]{
  reduced Planck constant
}

% 6. 工程、資訊與科技 Engineering, computing, and technology
\item[$O(\cdot)$]{
  asymptotic upper bound on a function's growth rate
}

% 7. 生物、農業與環境 Life science, agriculture, and environment
\item[$\mathrm{LD}_{50}$]{
  median lethal dose, the dose fatal to half of a test population
}

% 8. 醫學、健康與公衛 Medicine, health, and public health
\item[OR]{
  odds ratio
}

\end{denotation}
 那一行註解掉。
%
% 選用參數 [3cm] 是標籤欄寬度，預設 2.5cm。標籤比欄寬長時說明會被推開，
% 把數字調大即可。
%
% One example per broad field, covering the three kinds of label this
% environment handles: maths, typographic marks, and plain abbreviations. Keep
% the ones that fit your discipline and delete the rest. The optional [3cm] is
% the label column width; widen it if a label overruns.

\begin{denotation}[3cm]

% 1. 人文與文化 Humanities and culture
\item[$\dagger$]{
  obelus, marking a corrupt or doubtful passage in a critical edition
}

% 2. 社會、政治與公共政策 Social science, politics, and public policy
\item[GDP]{
  gross domestic product
}

% 3. 法律與制度 Law and institutions
\item[\S]{
  section sign, citing an article of a statute
}

% 4. 商業與管理 Business and management
\item[ROI]{
  return on investment
}

% 5. 基礎科學與數學 Basic science and mathematics
\item[$\hbar$]{
  reduced Planck constant
}

% 6. 工程、資訊與科技 Engineering, computing, and technology
\item[$O(\cdot)$]{
  asymptotic upper bound on a function's growth rate
}

% 7. 生物、農業與環境 Life science, agriculture, and environment
\item[$\mathrm{LD}_{50}$]{
  median lethal dose, the dose fatal to half of a test population
}

% 8. 醫學、健康與公衛 Medicine, health, and public health
\item[OR]{
  odds ratio
}

\end{denotation}
 那一行註解掉。
%
% 選用參數 [3cm] 是標籤欄寬度，預設 2.5cm。標籤比欄寬長時說明會被推開，
% 把數字調大即可。
%
% One example per broad field, covering the three kinds of label this
% environment handles: maths, typographic marks, and plain abbreviations. Keep
% the ones that fit your discipline and delete the rest. The optional [3cm] is
% the label column width; widen it if a label overruns.

\begin{denotation}[3cm]

% 1. 人文與文化 Humanities and culture
\item[$\dagger$]{
  obelus, marking a corrupt or doubtful passage in a critical edition
}

% 2. 社會、政治與公共政策 Social science, politics, and public policy
\item[GDP]{
  gross domestic product
}

% 3. 法律與制度 Law and institutions
\item[\S]{
  section sign, citing an article of a statute
}

% 4. 商業與管理 Business and management
\item[ROI]{
  return on investment
}

% 5. 基礎科學與數學 Basic science and mathematics
\item[$\hbar$]{
  reduced Planck constant
}

% 6. 工程、資訊與科技 Engineering, computing, and technology
\item[$O(\cdot)$]{
  asymptotic upper bound on a function's growth rate
}

% 7. 生物、農業與環境 Life science, agriculture, and environment
\item[$\mathrm{LD}_{50}$]{
  median lethal dose, the dose fatal to half of a test population
}

% 8. 醫學、健康與公衛 Medicine, health, and public health
\item[OR]{
  odds ratio
}

\end{denotation}
